\subsection{From Gaussian nilsequence estimates to the polynomial count}
\label{subsec:gaussian-gowers-bridge}

Fix the Gaussian ray field $K/\Q(i)$ and the conjugation-invariant ray
classes used in the preceding argument. Let $D$ be divisible by the
conductor of $K^{\rm ab}/\Q$. For each relevant union $C$ of classes,
let $\Lambda'_C$ be its rational prime weight and let $\mu_C(b)$ be the
conditional density in \eqref{eq:v5-mu}. In particular
$0\leq\mu_C(b)\leq1$. We use the following conclusion of the preceding
Gaussian norm, progression and factorization arguments, as assembled in
Corollary~\ref{cor:ga-fixed-centered-nilsequence}: for every fixed
multiple $V$ of $D$, every unit $b\pmod V$, every fixed bounded-complexity
nilsequence $F(g(n)\Gamma)$, and every $B>0$,
\begin{equation}\label{eq:gaussian-centered-nil-input}
\begin{aligned}
 \left|\mathbb E_{n\leq L}h_{C,V,b}(n)F(g(n)\Gamma)\right|
 &\ll_{V,B,\mathrm{complexity}}(\log L)^{-B},\\
 h_{C,V,b}(n)&=\frac{\varphi(V)}V
 \bigl(\Lambda'_C(Vn+b)-\mu_C(b)\Lambda'(Vn+b)\bigr).
\end{aligned}
\end{equation}
The bound is uniform in the polynomial orbit and functions with the stated
complexity bound. The field, its ray modulus and $V$ are fixed in this
statement. We do not assume any uniform rate as $V$ varies.

\paragraph{Published transference inputs.}
We use Green--Tao, \emph{Linear equations in primes}~\cite{greentao},
Proposition~6.4 and Appendix~D for the ordinary prime majorant,
Proposition~7.1 and Appendix~C for its generalized von Neumann bound,
and Lemma~4.4 with its subsequent normal-form reduction.
These are statements about the ordinary prime majorant and arbitrary
functions dominated by it; they do not assert an Artin-prime theorem.
We also use Dodos--Kanellopoulos~\cite{v5dk}, Theorem~5.1, in the
following form, checked in arXiv:1610.00487v4, pp.~12--13. The final
typeset statement has not been compared. For target interval norm $U^s[L]$, $s\geq2$, its hypothesis on a
nonnegative majorant on a prime cyclic group of size
$L'\in[CL,2CL]$, $C\geq20$, is
\begin{equation}\label{eq:gaussian-dk-correct-exponent}
 \|\nu-1\|_{U^{2s}(\Z/L'\Z)}\leq\eta.
\end{equation}
For a real function $|f|\leq\nu$ with $\|f\|_{U^s[L]}\geq\delta$,
the conclusion is a correlation bounded below with a bounded Lipschitz
linear nilsequence on an $(s-1)$-step nilmanifold. Its nilmanifold,
Lipschitz bound and positive correlation threshold depend only on
$s,C,\delta$. The interval norm is the normalized zero-extension norm.
In particular a target $U^{u+1}$ requires majorant norm $U^{2u+2}$.

\paragraph{An auxiliary cutoff, distinct from the polynomial cutoff.}
Put
\[
 V_z=D\prod_{\substack{p\leq z\\p\nmid D}}p.
\]
Here $z$ is an auxiliary cutoff used only to prove the next linear
replacement statement. It is not the cutoff $v$ defining $W,A$ in
\eqref{eq:v5-global-gowers}.
For every fixed integer $s\geq2$, one can choose $z=z(L)\to\infty$
sufficiently slowly that
\begin{equation}\label{eq:gaussian-refined-gowers}
 \max_{C,\ (b,V_z)=1}\|h_{C,V_z,b}\|_{U^s[L]}=o(1).
\end{equation}
The maximum is over the fixed finite list of ray-class unions and the
finitely many unit residues for the current $V_z$.

Here are all normalization and limit details. First choose the complexity
of the ordinary prime majorant larger than both the requirements of the
intended generalized von Neumann bound and
$\max(4^s,4s)$. Choose its ambient constant $C\geq20$ large enough for
both published majorant and generalized von Neumann statements. If
$\lambda_{V,b}(n)=(\varphi(V)/V)\Lambda'(Vn+b)$, then
$|h_{C,V,b}|\leq\lambda_{V,b}$. The majorant simultaneously dominates
a fixed number of such ordinary weights on $[L^{3/5},L]$, with a
multiplicative constant depending only on the chosen complexities.
Consequently there is $\alpha>0$, depending only on the fixed
complexities and the fixed factor in the modulus, such that
$\alpha 1_{[L^{3/5},L]}h_{C,V_z,b}$ is dominated in absolute value by
$\nu$ on $[L]$. The same $\alpha$ works for all cutoffs and unit residues.
In applying the inverse theorem to a lower bound $\delta$ for the
unscaled norm, use the lower bound $\alpha\delta/2$ for this truncated,
scaled function, once the truncation error is at most $\delta/2$.

For sufficiently large $z$, $V_z=E\prod_{p\leq z}p$ with $E$ fixed.
The ordinary divisor-sum construction works with this $V_z$: at primes
dividing $E$, the unit residue excludes every divisor term containing
that prime, and at all other primes its local factors are unchanged.
The same observation applies to the correlation upper bounds in
Appendix~D. Thus the stated pseudorandomness conditions remain valid,
uniformly in unit residues, with constants allowed to depend on $E$.
Expand the centered $(2s)$-cube of $\nu-1$. It contains at most $4^s$
distinct forms, in $2s+1$ variables, with coefficients zero or one.
Every subproduct has mean $1+o(1)$ by the linear-forms condition;
inclusion--exclusion gives \eqref{eq:gaussian-dk-correct-exponent} with
a left side tending to zero. This estimate is on the full cyclic group.

The removed initial interval contributes $o(1)$ to the interval $U^s$
norm. Indeed the weights are $O(\log L)$ if $V_z=L^{o(1)}$, and a cube
with its initial vertex in this interval has relative count
$O(L^{-2/5})$. The $2^s$th power of the norm of the removed function is
at most $O((\log L)^{2^s}L^{-2/5})$. Its correlation with a bounded
nilsequence is also $o(1)$.

If \eqref{eq:gaussian-refined-gowers} fails by some fixed $\delta>0$,
Theorem~5.1 gives a fixed-complexity nilsequence correlation bounded
away from zero, contradicting \eqref{eq:gaussian-centered-nil-input}.
There is no illicit varying-modulus use of that input. For each integer
$z$ and each of the finitely many residues modulo $V_z$, first take a
size threshold for the required nilsequence estimates. Diagonalize
these thresholds also over $\delta=1,1/2,1/3,\ldots$ and over the finite
list of norm orders needed. Impose in addition
$z(L)\leq\log\log\log L$ and all size thresholds of the majorant.
At stage $j$, choose the cutoff $z_j$ large enough for the limiting
majorant errors to be below half the inverse-theorem tolerances needed
for the first $j$ tests; only then choose their size threshold $L_j$.
The resulting increasing step function tends to infinity and proves
the assertion for all sufficiently large $L$. The fixed smooth metric
in the inverse theorem is comparable to a fixed Mal'cev metric, and a
linear orbit with arbitrary base point is a polynomial orbit, so the
nilsequence input applies to exactly the supplied correlations.
When this construction is used at the refined scale $L\asymp N/V_{z_j}$
below, choose the physical threshold $N_j$ so large that every constant
multiple of $N/V_{z_j}$ used by the normal-form reduction is at least
$L_j$ for $N\geq N_j$. Impose the same growth constraints there.
Thus the cutoff is constant on each physical-scale stage; no implicit
equation $z=z(N/V_z)$ is being assumed.

\begin{lemma}[Fixed-system linear replacement]
\label{lem:gaussian-linear-replacement}
Let $\Psi=(\psi_1,\ldots,\psi_t):\Z^d\to\Z^t$ have fixed nonzero integer
linear parts of finite complexity. Allow its constant terms to vary
with $N$, subject to $\|\Psi\|_N\leq L_0$ for a fixed $L_0$.
Uniformly over convex regions $\mathcal K\subseteq[-N,N]^d$ with
$\Psi(\mathcal K)\subseteq[1,L_0N]^t$,
\begin{equation}\label{eq:gaussian-linear-replacement}
 \sum_{\mathbf n\in\mathcal K\cap\Z^d}
 \left(\prod_{i=1}^t\Lambda'_{C_i}(\psi_i(\mathbf n))
 -\prod_{i=1}^t\mu_{C_i}(\psi_i(\mathbf n)\bmod D)
                         \Lambda'(\psi_i(\mathbf n))\right)=o(N^d).
\end{equation}
The same conclusion holds for a mixed product in which any selected
factors have already been replaced by their ordinary-prime models.
\end{lemma}
\begin{proof}
We can discard points where some $\psi_i\leq N^{9/10}$.
For a nonconstant fixed linear form their number is
$O_{\Psi,L_0}(N^{d-1}(N^{9/10}+1))$; the product of weights is
$O((\log N)^t)$. This is $o(N^d)$, uniformly in the constant terms.
Apply the auxiliary $V_z$-trick to the variables:
$\mathbf n=V_z\mathbf m+\mathbf r$.
If one output residue is not a unit modulo $V_z$, its prime weight
vanishes on the remaining range, since the output is larger than
every prime dividing $V_z$. For an admissible residue vector write
\[
 \psi_i(V_z\mathbf m+\mathbf r)
       =V_z\widetilde\psi_{i,\mathbf r}(\mathbf m)+b_{i,\mathbf r},
       \qquad1\leq b_{i,\mathbf r}\leq V_z.
\]
The new linear coefficients are the original ones. Its normalized
constant terms remain bounded at scale $N/V_z$. Because $D\mid V_z$,
the factor $\mu_{C_i}$ is constant on this residue vector.

Telescope the difference of the two products. After normalization by
$(\varphi(V_z)/V_z)^t$, each term has one factor $h_{C_i,V_z,b_i}$;
all other factors are nonnegative and at most
$\lambda_{V_z,b_j}$. Introduce a common ordinary-prime majorant for
these at most $t$ weights and truncate as above. Multiply every
factor by a common fixed $\alpha_0>0$ so that each has absolute value
at most that majorant, and apply the normal-form extension of
Green--Tao Lemma~4.4 and Proposition~7.1. Here $\alpha_0$ depends
only on the fixed system, the chosen complexities and the fixed
conductor, not on $z$ or the residue vector. The small factor has
the required norm by \eqref{eq:gaussian-refined-gowers}. Dividing
the resulting estimate by $\alpha_0^t$ changes only a fixed constant.
The normal-form extension adds only a fixed number of variables with
fixed coefficients; its convex-region reduction is the one following
Lemma~4.4, and divides out their box volume. The error on each residue
vector is therefore $o((N/V_z)^d)$, uniformly in the vector and convex
region. The discarded initial segment stays outside the retained
output range since $V_z=N^{o(1)}$.

One can make this reduction explicit. For each pair $i\ne j$ choose
an integer vector $v_{ij}$ with $\dot\psi_j(v_{ij})=0$ and
$\dot\psi_i(v_{ij})\ne0$; fixed nonproportional linear parts ensure
such vectors exist with fixed size. Replace $\mathbf m$ by
$\mathbf x+\sum_{i\ne j}y_{ij}v_{ij}$ and allow each new $y_{ij}$
to range over an integer interval of length comparable to $N/V_z$.
The block $y_{ij}$ for fixed $i$ certifies normal form. The region
defined by the original convex constraint on this sum and the new box
is convex. The integer change $(\mathbf x,\mathbf y)\mapsto
(\mathbf x+\sum y_{ij}v_{ij},\mathbf y)$ is bijective, so its sum
is exactly the original sum times the number of points in the new box.
Division by that number gives the asserted error. If there is only one
form, choose an integer vector $v$ with $\dot\psi_1(v)\ne0$ and replace
$\mathbf m$ by $\mathbf x+yv$. The same box and bijection argument
gives the reduction in this case.

No growing factor is lost when summing residue vectors. If
$\mathcal R_z$ is the set of admissible vectors, then
\begin{equation}\label{eq:gaussian-admissibility-mass}
 \frac{|\mathcal R_z|}{V_z^d}
       \left(\frac{V_z}{\varphi(V_z)}\right)^t
   =\prod_{p\mid V_z}\beta_p(\Psi)=O_{\Psi,L_0}(1).
\end{equation}
Exclude the finitely many primes dividing a chosen nonzero
coefficient of each linear part, and a chosen nonzero two-by-two
minor for each pair of linear parts. At every remaining prime each
output hyperplane has probability $p^{-1}$, and any two meet with
probability $p^{-2}$. Hence $\beta_p=1+O_t(p^{-2})$, uniformly in
the constant terms. This also covers $t=1$, when there are no pair
minors and the good-prime factor is exactly one.
The finitely many remaining factors are bounded by $(p/(p-1))^t$.
The Chinese remainder theorem proves the equality, also for the fixed
prime powers dividing $D$. Multiplying the per-residue error by this
bounded mass gives $o(N^d)$. The proof works unchanged for the mixed
products. All auxiliary thresholds can be chosen uniformly over bounded
constant terms and convex regions, because the published majorant and
generalized von Neumann conclusions have this uniformity.
\end{proof}

\paragraph{The exact local cube needed in the appendix.}
Now fix the appendix cutoff $v$, and put $W,A,n,R$ as in its proof.
This is a new limiting argument; the auxiliary cutoff in the proof of
Lemma~\ref{lem:gaussian-linear-replacement} has already been eliminated
from that lemma's statement. Let
\[
 f_{C,b}(x)=\frac{\varphi(W)}W\Lambda'_C(Wx+b)1_{R\leq x\leq n},
 \qquad g_{C,b}=f_{C,b}-\mu_C(b).
\]
For fixed $a\geq2$ and $C_1<C_2<C_3$, uniformly over
$q\leq A^{C_1}$ and $nA^{-C_3}\leq J\leq nA^{-C_2}$,
\begin{equation}\label{eq:gaussian-global-gowers}
 \|g_{C,b}\|_{U^a_{q[J]}}^{2^a}
       \ll_{a,C_1,C_2,C_3,\rho,D}\frac1{\log v}+o_v(1).
\end{equation}
The norm uses the difference multiset $q(h-h')$, $h,h'\in[J]$, on
$\Z/n\Z$.

To prove this, remove the cutoff $x<R$ at cost
$O((\log n)^{2^a}R/n)=o_v(1)$. Expand the centered cube and partition
its domain according to every wrap-around integer $m_\omega$.
The relevant affine forms are precisely
\[
 W\left(x+q\sum_{r=1}^a\omega_r(h_r-h'_r)-m_\omega n\right)+b,
 \qquad\omega\in\{0,1\}^a.
\]
They have pairwise nonproportional linear parts. Each region is convex;
there are $O_a(1)$ regions because $qJ/n<1$. For fixed $v$, the number
of possible $q$ is finite and $J/n$ stays in a fixed positive compact
interval. Apply Lemma~\ref{lem:gaussian-linear-replacement} to every
subproduct and region. Its error is $o_v(n^{2a+1})$ uniformly in $J$,
and division by $nJ^{2a}$ gives $o_v(1)$. Since every output is $b$
modulo $D$, all its conditional densities are the same $\mu_C(b)$.
Thus the centered Artin cube equals $\mu_C(b)^{2^a}$ times the
centered ordinary-prime cube, plus $o_v(1)$.

The ordinary cube is evaluated exactly as in Tao--Ziegler~\cite{v5tz},
Proposition~6, published pages 39--40, using the ordinary-prime linear
correlation theorem, for example the published formulation in
Bienvenu~\cite{v6bienvenu}, Theorem~1.1, pages 79--80.
The factors at $p\mid W$ cancel the normalization.
For $p\nmid Wq$, pairwise independent gradients give local factor
$1+O_a(p^{-2})$, uniformly in the wrap constants. For $p\mid q$,
$p>v$, the factor is $1+O_a(p^{-1})$. Since
$\log q=O_{C_1,\rho,D}(v)$, there are $O(v/\log v)$ such prime divisors.
Their product is $1+O(1/\log v)$, and the other large-prime factors
give $1+O(1/v)$. Centering cancels the common main term. Prime powers,
if the ordinary theorem is stated for $\Lambda$, contribute $o_v(1)$
by their $O(\sqrt n\log^2 n)$ total weight and the fixed linear-fiber
bound. This proves \eqref{eq:gaussian-global-gowers}, and hence supplies
\eqref{eq:v5-global-gowers} without an Artin linear-correlations theorem.

\paragraph{Completion by the published polynomial transfer.}
The remainder of the proof of Proposition~\ref{prop:v5-artin-patterns},
starting with its arbitrary-function dense-model step, now applies
unchanged to these fixed Gaussian ray fields. Specifically use the
published Tao--Ziegler~\cite{v5tz} Theorem~11 and Corollary~3, then Theorem~9 and
Theorem~13. The model for an arbitrary $f\leq\lambda_{W,b}$ is used;
one does not assert that an Artin weight has model one. Equation
\eqref{eq:gaussian-global-gowers} shows that its bounded model, minus
$\mu_C(b)$, has small global norms. The model is bounded by a fixed
constant depending on the chosen sieve exponent. Divide its difference
from $\mu_C(b)$ by a fixed bound for that difference before applying
Theorem~9, whose hypothesis is a bound by one. After restoring this
fixed factor, concatenation and model approximation give the required
local norm bounds. None of these scaling constants depends on $v$.

Keep the original parameter order: local box order and coefficient
bound first, then $B_0>C+2$, then the global order $a$, and finally
small $\rho$. Only then let $v$ increase. The quantities
$M_0=\lfloor T/W\rfloor$, $M_1=\lfloor M_0/A\rfloor$ and
$J=\lfloor A^{-2B_0a}M_1^2\rfloor$ satisfy the size hypotheses recorded
in the appendix when $T/\sqrt N\geq4\sqrt W/A$.
For each fixed $v$ the estimates are uniform for $T$ in that allowed
range. The final assertion $T=o(\sqrt N)$ is imposed only along the
subsequent diagonal $v=v(N)\to\infty$, not while $v$ is frozen.

After $u=Wx+b$, $w=Wm+c$, the polynomials are
$c_jWm^2+2c_jcm+(b+c_jc^2-b_j)/W$, with distinct leading coefficients.
The bounded-model transfer kills every term containing a centered
factor. Output-boundary strips have width $O(T^2)=o(N)$.
For completeness, dominate them first by strips of width $\epsilon N$.
Here is the needed localization of the ordinary majorant estimate.
This argument also covers negative coefficients $c_j$. For the true
integer count use an enlarged period $n_* = \lceil B N/W\rceil$, where
$B>2+\max_j|c_j|$ is fixed. We may assume $T\leq\sqrt N$.
Nonpositive outputs have prime weight zero, and every positive output
is at most $(1+\max_j|c_j|)N$. Discard also tuples for which some
positive output is at most $W\lceil n_*^\rho\rceil+W$. For each $w$
there are $O_k(W n_*^\rho+W)$ such $u$, so their total weight is
$O_k((W n_*^\rho+W)T(\log N)^k)=o_v(NT)$, uniformly in the
allowed $T$ range. All remaining true outputs have their ordinary
prime majorants on $[n_*]$ without wrap-around. Use the same
localization calculation on the original period $n$ to bound the
artificial wrapped terms of the cyclic count. The parameters $n_*$
and $n$ differ by a fixed factor, so both applications satisfy the
same polynomial-forms hypotheses. In the following calculation $n$
denotes either chosen ambient period.
Write $\nu_0(x)$ for the majorant factor with constant shift if there
is one, and set $\nu_0=1$ otherwise. There is at most one such factor.
Put $F(x)=\mathbb E_m\prod_i\nu_i(x+P_i(m))$, omitting that factor.
Tao--Ziegler, Theorem~15 in Appendix~A (published page~57), applied to
$\nu_0F$, $\nu_0F^2$ and $\nu_0$, gives
\[
 \mathbb E_x\nu_0(x)(F(x)-1)^2=o_{v\to\infty}(1)+o_v(1).
\]
In the square use independent variables $m,m'$; the differences
between all its nonconstant shift polynomials are nonconstant,
including those from distinct copies. Weighted Cauchy--Schwarz now
shows, for every interval $I$ in the $x$ variable,
\[
 \mathbb E_x1_I(x)\nu_0(x)F(x)
 =\mathbb E_x1_I(x)\nu_0(x)+o_{v\to\infty}(1)+o_v(1).
\]
For fixed relative interval length the one-dimensional divisor-sum
mean estimate gives
$\mathbb E_x1_I\nu_0=|I|/n+o_{v\to\infty}(1)+o_v(1)$;
it follows either directly from the ordinary majorant construction
or from its linear-forms estimate on that fixed interval.
The normalized prime weights are bounded by fixed multiples of
their sieve majorants, with constants depending on the fixed sieve
exponent. Thus a strip of width $\epsilon N$ has total weight
$[O(\epsilon)+o_{v\to\infty}(1)+o_{v,\epsilon}(1)]NT$, with that
fixed domination factor included in the implied constants.
The errors here have the same iterated order as the rest of the
polynomial transfer. Then take $\epsilon\downarrow0$ as described
below. This avoids a spurious factor $(\log N)^k$ from a pointwise
bound. The same positive majorant treats interval endpoints.

Choose increasing cutoffs $v_j$ for which the global/local errors and
$4\sqrt{W(v_j)}/A(v_j)$ tend to zero, and then choose size thresholds
valid for all the finitely many residues, step values and transfer
parameters at that cutoff.
At stage $j$ impose the boundary estimates for every
$\epsilon\in\{1,1/2,\ldots,1/j\}$, preserving all earlier tests.
Let $v(N)=v_j$ between successive thresholds and take
$\omega(N)\geq4\sqrt{W(v(N))}/A(v(N))$, tending to zero sufficiently
slowly. This gives the uniform assertion for
$\omega(N)\sqrt N\leq T=o(\sqrt N)$.
Indeed, for an arbitrary such $T$, first fix $k$; eventually
$T^2/N<1/k$ and the estimate for the strip of width $N/k$ applies.
Take the upper limit as $N\to\infty$ and then $k\to\infty$.
There is no requirement that $T^2/N<1/j(N)$ along the chosen stages.
The main term is precisely \eqref{eq:v5-artin-series}; the admissible
residue mass is bounded, so summation does not amplify errors.
The $u\in[N/2,N]$ version follows by using that convex interval directly,
or by subtraction with thresholds enlarged for $N/2$.

This bridge uses the separately established Gaussian nilsequence input
and published ordinary-prime transference. It uses no statement or
estimate from Frei--Tinkov\'a as a theorem.


